\section{Séquences, replis et régimes}

\subsection{Combien de pertes d'affilée faut-il supporter}

\begin{figure}[H]
\centering
\includegraphics[width=\linewidth]{usdjpy_trades_streaks}
\caption{Séries de gains et de pertes consécutifs, dans l'ordre des clôtures.
Les séries de trois trades ou plus portent leur longueur.}
\end{figure}

À première vue, ce candidat est bien plus confortable que l'or~: la plus longue
série perdante ne compte que \textbf{trois} trades, contre dix sur la sleeve or,
et les alternances gain/perte se succèdent sans jamais former de longue
traversée du désert. Un investisseur qui compterait ses pertes consécutives
tiendrait sans difficulté.

Ce serait la mauvaise métrique.

\begin{warningbox}
La plus longue série perdante --- trois clôtures, du 21~mars au 3~mai~2023 ---
coûte \textbf{$-$96\,940\,\$}, soit \metric{33,8\,\%} du sous-compte en
\textbf{six semaines}. Ce n'est pas la \emph{longueur} des séries qui fait le
risque de ce candidat, c'est leur \emph{profondeur}~: trois erreurs consécutives
suffisent, parce que chacune est portée par un levier proche du plafond. La
deuxième plus mauvaise séquence, en septembre~2025, retire encore
$-$45\,343\,\$ en trois semaines.
\end{warningbox}

C'est une différence de nature avec la sleeve or, et elle a une conséquence
opérationnelle~: la douleur n'y est pas étalée sur huit mois d'ennui, elle est
concentrée sur quelques semaines de pertes rapides. Les deux exigent la même
discipline, mais elles ne se ressemblent pas et ne se rassurent pas de la même
manière.

\subsection{Le repli}

\begin{figure}[H]
\centering
\includegraphics[width=\linewidth]{usdjpy_trades_equity_dd}
\caption{Balance du sous-compte USD/JPY et repli associé. Les trois plus fortes
descentes sont annotées.}
\end{figure}

Le repli maximal mesuré sur cette balance atteint \metric{33,8\,\%}, et il
coïncide au dollar près avec le repli de balance rapporté par le simulateur
(96\,940,43\,\$, soit 33,79\,\%). Deux précautions de lecture s'imposent malgré
tout.

Premièrement, la balance ne bouge \emph{qu'à la clôture} d'une position. Entre
deux clôtures, le capital réellement exposé fluctue sans que cette courbe le
montre. Le simulateur, qui mesure le repli tick par tick, rapporte pour le même
run un repli d'équité de \textbf{45,2\,\%} --- onze points de plus. Le chiffre
de 33,8\,\% est donc une borne \textbf{inférieure} du risque subi, et c'est
45,2\,\% qu'il faut retenir.

Deuxièmement, ce repli est celui d'une sleeve isolée jouée à pleine allocation
avec un facteur de risque de 1,0. Il n'a pas de sens transposé tel quel à un
portefeuille~: ni l'échelle, ni la pondération, ni la mutualisation avec les
autres moteurs ne sont représentées ici.

\subsection{Par année}

\begin{figure}[H]
\centering
\includegraphics[width=\linewidth]{usdjpy_trades_yearly}
\caption{Résultat net par année de clôture, avec le nombre de trades gagnants
sur le total.}
\end{figure}

La ventilation annuelle est le meilleur résumé visuel du problème de ce
candidat~: une barre, 2022, à $+$188\,374\,\$ pour \emph{un} trade~; et cinq
autres années qui, ensemble, retirent 11\,551\,\$. 2025 est la plus mauvaise,
avec neuf trades pour $-$79\,320\,\$ et trois gagnants seulement~; 2023 suit,
avec dix trades pour $-$26\,247\,\$.

Le découpage est fait par \emph{année de clôture}, ce qui déplace du résultat
d'une barre à l'autre~: le trade de 2021 est crédité à 2022, celui d'octobre
2025 à 2026. La barre 2026 ($+$42\,695\,\$) est particulièrement trompeuse~:
elle contient une position ouverte en octobre~2025 et une autre liquidée par
le testeur en fin de fenêtre. Le chapitre~\ref{sec:limites} y revient, parce
que cette même ligne sert de lecture hors échantillon.

Il faut enfin résister à la tentation d'y lire un régime. Avec six années
civiles, trente-cinq trades et une seule position qui porte tout, chaque barre
annuelle repose sur une poignée d'observations~; ces chiffres décrivent ce qui
s'est produit, ils ne prédisent rien.
